% ---------------------------------------------------------------------------
%  How this book is arranged
% ---------------------------------------------------------------------------
\chapter*{How This Book Is Arranged}
\addcontentsline{toc}{chapter}{How This Book Is Arranged}
\markboth{How This Book Is Arranged}{}
\thispagestyle{plain}

This is the whole of HS248AT in one document: every question recorded for the
course --- \textbf{150 in all} --- answered, plus every paper those questions
came from, reproduced as it was set.

\vspace{2pt}
It is in three parts, because there are three different ways you will want to
use it.

\begin{abul}
\item \textbf{Part I --- the descriptive question bank.} Every long-answer
      question ever set, arranged \emph{unit-wise} against the syllabus, each
      with a model answer. This is the part you revise from. Reading it
      straight through covers the course; \textbf{74 questions}.
\item \textbf{Part II --- the objective question bank.} Every multiple-choice
      question, fill-in-the-blank and one-mark short answer, with the answer
      beside it. \textbf{76 items}, in three tables --- 36 multiple-choice,
      35 fill-in-the-blank, 5 short answers.
\item \textbf{Part III --- the question papers.} All nine papers reproduced in
      full and in chronological order, with the date, assessment type,
      duration, maximum marks and the printed instructions. Every question
      carries a pointer to its answer in Part~I or Part~II, so a whole past
      paper can be worked end to end and then marked.
\end{abul}

The two arrangements are deliberate. Part~I is arranged by topic because
that is how you learn; Part~III is arranged by paper because that is how you
are examined, and a paper's shape --- how many marks Part~A carries, which
questions offer a choice --- is itself worth knowing before you sit down.

\subsection*{Where the answers come from}
\addcontentsline{toc}{section}{Where the answers come from}

Every answer in this book is written \textbf{strictly from the prescribed
material and nothing else}: the three unit notes for this course (\emph{A
Foundation Course in Human Values and Professional Ethics}, Chapters 1--10),
and the official RVCE \emph{Scheme \& Solution} documents for CIE~2 (2022--23)
and Quiz-3\,/\,CIE~3 (2023--24). Where a scheme gives the model answer, that
wording is followed in substance.

Nothing has been imported from outside those sources. Where a past-paper
question uses a term the notes do not define in that form, the answer says so
and answers from the nearest thing the notes do define, rather than filling
the gap from elsewhere.

\begin{note}
\textbf{\sffamily Two warnings about circulating answer-sets.}

\textbf{(1) Justice.} The Unit~II answer-set circulating as ``unit-2 uhv
question bank answers'' defines justice using the legal categories
\emph{distributive\,/\,procedural\,/\,retributive\,/\,restorative}. That
contradicts the official RVCE scheme, which defines the four elements of
justice as \emph{recognition of values, fulfilment, evaluation,} and
\emph{mutual happiness ensured}. This book follows the official scheme and the
prescribed text throughout --- write those in the exam.

\textbf{(2) Chapters 11--13.} Two SEE questions (\emph{holistic technology};
\emph{ethical human conduct in terms of values, policy and character}) belong
to Chapters 11--13, which are not part of the Unit I--III material supplied.
They are listed in the paper index but not answered here, since no source was
provided for them. That omission is flagged where it occurs rather than
papered over.
\end{note}

\subsection*{Reading the tags}
\addcontentsline{toc}{section}{Reading the tags}

Each question in Part~I carries a tag block naming the papers it has actually
appeared in, and the marks the source paper allotted:

\vspace{4pt}
\begin{keybox}
\textbf{Q32.}\enspace What do you understand by trust? Differentiate between
intention and competence.\tg{CIE 2}\tg{SEE}\mmk{5}

\vspace{5pt}
{\footnotesize This question was set in a CIE~2 paper and in an SEE, and
carried 5 marks. A question tagged with three or four papers is the single
best use of revision time --- Appendix~B lists them all.}
\end{keybox}

\vspace{4pt}
\tg{Scheme} means the answer follows the official RVCE scheme verbatim in
substance. An untagged question has not yet appeared in a paper we hold; it is
in the book because the departmental question bank or the notes' own review
questions set it, and it covers syllabus the papers have so far skirted.

In Part~III every question additionally carries the three codes its paper
printed beside it --- \textbf{M} (marks), \textbf{CO} (course outcome, code
only) and \textbf{BT} (Bloom's level, \textbf{L1} remember, \textbf{L2}
understand, \textbf{L3} apply, \textbf{L4} analyse). No paper in this
collection sets anything above L4.

\subsection*{The nine papers}
\addcontentsline{toc}{section}{The nine papers}

\renewcommand{\arraystretch}{1.35}
\arrayrulecolor{rulegrey}
\begin{tabular}{@{}P{46mm} P{32mm} C{20mm} C{20mm} C{28mm}@{}}
\rowcolor{slate}
\hdr{Assessment} & \hdr{Date} & \hdr{Max M} & \hdr{Time} & \hdr{Reproduced} \\
\hline
SEE (21HSU48) & October 2023 & 50 & 2 h & Index only \\ \hline
CIE 1 & 22 June 2024 & 25 & 60 min & Index only \\ \hline
CIE 2 & 24 July 2024 & 5 + 25 & 60 min & Index only \\ \hline
Improvement CIE\,/\,CIE 3 & 28 August 2024 & 25 & 60 min &
  \textbf{Part III} \\ \hline
SEE & Sept/Oct 2024 & 50 & 2 h & Index only \\ \hline
CIE-II & 7 May 2025 & 25 + 5 & 60 min & \textbf{Part III} \\ \hline
Improvement CIE & June 2025 & 30 & 60 min & \textbf{Part III} \\ \hline
\textbf{SEE} Regular\,/\,Suppl. & June/July 2025 & 50 & 2 h &
  \textbf{Part III} \\ \hline
Improvement CIE & 18 June 2026 & 05 + 25 & 10 + 60 min &
  \textbf{Part III} \\ \hline
\end{tabular}

\vspace{6pt}
Five papers are reproduced question-for-question in Part~III because scans of
the full papers were available. For the other four only the question list
survives, so they appear in the Appendix~A index --- but every one of their
questions is answered in Parts~I and~II just the same. Nothing has been
dropped for want of a scan.

\subsection*{How long an answer should be}
\addcontentsline{toc}{section}{How long an answer should be}

Answers here are written to the length the marks buy. The valuation scheme for
this course awards roughly one mark per distinct, correctly explained point,
so a 5-mark answer carries five scoring points, a 6-mark answer six, an 8-mark
answer eight. A 4-mark definition question is two to four sentences.

In the examination you are not rewarded for volume. You are rewarded for
hitting the number of distinct points the scheme expects, each named and then
explained in a line or two. Where an answer below runs longer than you could
write in the time available, treat the \textbf{numbered points as the skeleton
to reproduce} and the prose as the explanation to compress.

\vspace{4pt}
Appendix~C lists, unit by unit, what to revise first if you have run out of
time.
